\documentclass[11pt,a4paper]{article}



\usepackage{fontspec}
\usepackage{polyglossia}
\setmainlanguage{english}
\setotherlanguage{russian}
\setmainfont{FreeSerif}
\newfontfamily\cyrillicfont{FreeSerif}


\usepackage[margin=2.5cm]{geometry}
\usepackage{amsmath,amssymb}
\usepackage{booktabs}
\usepackage{array}
\usepackage{multirow}
\usepackage{hyperref}
\usepackage{xcolor}
\usepackage{graphicx}

\usepackage{enumitem}
\usepackage{caption}
\usepackage{float}
\usepackage{setspace}
\usepackage{parskip}

\hypersetup{
    colorlinks=true,
    linkcolor=blue!60!black,
    citecolor=blue!60!black,
    urlcolor=blue!60!black
}

\setlength{\parindent}{0pt}
\setlength{\parskip}{6pt}

% ── Title ──────────────────────────────────────────────────────────────────
\title{\textbf{PhenoSeq: Bilingual Clinical Phenotype Extraction\\
and HPO Linking for Rare Disease Diagnosis}}
\author{Constantine Rosljakov}
\date{2026}

\begin{document}

\maketitle
\thispagestyle{empty}

% ── Abstract ───────────────────────────────────────────────────────────────
\begin{abstract}
We present \textbf{PhenoSeq}, an end-to-end system for automated extraction of
clinical phenotypes from free-text medical descriptions in English and Russian,
mapping them to Human Phenotype Ontology (HPO) terms, and generating
Orphanet-based differential diagnoses with Phenopackets~v2 export.  The system
combines a biomedical named entity recognition (NER) model
(\texttt{d4data/biomedical-ner-all}) with dictionary matching for English, and
a four-layer hybrid pipeline (dictionary, NER, cross-lingual retrieval,
rule-based) for Russian.  HPO linking uses multilingual sentence embeddings
(\texttt{paraphrase-multilingual-MiniLM-L12-v2}, 384-dim, 43{,}596-row
synonym-expanded index) with hierarchy-adjusted re-ranking and optional
cross-encoder reranking.  A statistical DNA conservation analysis module using
GC content, 3-mer Shannon entropy, and seven regulatory motifs complements the
phenotype extraction pipeline.

Evaluated on a 24-case bilingual test set (47~HPO annotations), the system
achieves $\text{Hit@3} = 0.830$ overall (English formal: 0.789--1.000, Russian
formal: 1.000, Russian informal: 0.333) with Wu--Palmer ontological similarity
$\overline{\mathrm{WP}}_3 = 0.872$.  The English NER component achieves
$F_1 = 0.713$ on a 30-sentence annotated corpus; the Russian NER component
achieves $F_1 = 0.765$ on an 8-sentence evaluation set.

\noindent\textit{Code:} \url{https://github.com/your-username/PhenoSeq}
\end{abstract}

\newpage
\tableofcontents
\newpage

% ═══════════════════════════════════════════════════════════════════════════
\section{Introduction}
% ═══════════════════════════════════════════════════════════════════════════

Rare disease diagnosis is a global clinical challenge.  Patients wait an
average of 4--6 years before receiving a correct diagnosis
\cite{nguengang2020}, partly because phenotype descriptions in medical records
are unstructured free text that cannot be automatically queried against disease
databases.  The Human Phenotype Ontology (HPO) \cite{kohler2021} provides a
standardised vocabulary of over 19{,}000 phenotypic abnormalities, and tools
such as Phenomizer \cite{kohler2009} and LIRICAL \cite{robinson2020} can rank
rare diseases given a set of HPO terms.  The bottleneck is the \emph{extraction
step}: transforming a clinician's narrative into structured HPO codes.

This problem is complicated by the multilingual nature of clinical practice.
Existing NLP tools for phenotype extraction are almost exclusively trained on
English text \cite{nevéol2018}, leaving Russian-speaking clinicians without
automated support despite a patient population of over 140~million.

PhenoSeq addresses both challenges by providing:
\begin{itemize}[nosep]
  \item A bilingual NER pipeline: English (NER + dictionary) and Russian
        (four-layer hybrid: dictionary $\to$ NER $\to$ cross-lingual retrieval
        $\to$ rule-based).
  \item Multilingual HPO linking via sentence transformers with synonym
        expansion and hierarchy-adjusted re-ranking.
  \item Orphanet differential diagnosis suggestions.
  \item Phenopackets~v2 export \cite{jacobsen2022}.
  \item \textbf{ConservSeq}: a statistical DNA conservation module that links
        patient symptoms to candidate disease genes and identifies functionally
        important DNA regions.  When clinicians receive sequencing results
        containing hundreds of mutations, ConservSeq helps them decide which
        mutations to investigate first --- those located in evolutionarily
        conserved regions are more likely to be pathogenic.
\end{itemize}

\subsection{Research Questions}

\begin{description}[nosep]
  \item[RQ1.] Can a combined NER + semantic linking pipeline achieve
        $\text{Hit@3} \geq 0.85$ on formal clinical text in both English and
        Russian?
  \item[RQ2.] Does synonym-expanded indexing and hierarchy-adjusted re-ranking
        improve MRR over a vanilla cosine-similarity baseline?
  \item[RQ3.] Can statistical DNA conservation signals (GC content, regulatory
        motifs, Shannon entropy) produce region labels that correlate with
        published phastCons scores?
\end{description}

% ═══════════════════════════════════════════════════════════════════════════
\section{Related Work}
% ═══════════════════════════════════════════════════════════════════════════

\subsection{Clinical NLP and Phenotyping}

\citeauthor{kohler2021} describe the HPO structure (19{,}000+ terms) and its
integration with OMIM, Orphanet, and ClinVar \cite{kohler2021}.
\citeauthor{robinson2010} define the IS-A hierarchy and semantic similarity
for phenotypic terms \cite{robinson2010}.  \citeauthor{jacobsen2022} define
Phenopackets~v2---the interchange format used for our export module
\cite{jacobsen2022}.

\citeauthor{groza2015} published the first systematic recall/precision
benchmark for HPO extraction tools, finding that recall drops below 60\% on
informal clinical language \cite{groza2015}.  \citeauthor{nevéol2018} surveyed
clinical NLP for languages other than English, showing that fewer than 5\% of
clinical NLP publications address non-English text \cite{nevéol2018}---directly
motivating the Russian NER component of this work.

\subsection{NLP Models for Biomedical Entity Recognition}

\citeauthor{lee2020} introduced BioBERT, the foundation model underlying most
modern biomedical NER systems \cite{lee2020}.  \citeauthor{huang2022} developed
PhenoBERT, fine-tuning BioBERT specifically on HPO-annotated clinical texts
\cite{huang2022}.  \citeauthor{neumann2019} released ScispaCy, a
CPU-efficient spaCy-based pipeline for biomedical NLP \cite{neumann2019}.

For English NER, PhenoSeq uses \texttt{d4data/biomedical-ner-all}, a general
biomedical NER model, combined with a comprehensive dictionary of 42{,}355 HPO
term names and synonyms for maximum recall.

\subsection{Russian Clinical NLP}

\citeauthor{miftah2017} established the RuMedNER corpus---the first Russian
biomedical NER benchmark \cite{miftah2017}.  \citeauthor{louka2021} released
NEREL-BIO, a dataset of 700 PubMed abstracts with disease/symptom annotations
in Russian and English \cite{louka2021}.  For Russian NER, PhenoSeq uses
\texttt{graviada/labse-ner-runne-ru}, supplemented with dictionary matching
(107 Russian HPO terms with full case-form inflection), cross-lingual LaBSE
retrieval, and rule-based extraction.

\subsection{Ontology-Aware Ranking}

\citeauthor{xiong2021} proposed ANCE for dense retrieval with hard negatives
\cite{xiong2021}.  \citeauthor{nogueira2020} introduced cross-encoder
reranking, which PhenoSeq implements as an optional post-retrieval step
\cite{nogueira2020}.  \citeauthor{kohler2009} defined ResnikSim and Lin
similarity---precursors to the Wu--Palmer metric used in our evaluation
\cite{kohler2009}.

\subsection{DNA Conservation}

\citeauthor{pollard2010} defined phastCons and phyloP scores, which serve as
ground truth for validating our statistical conservation method
\cite{pollard2010}.  \citeauthor{alipanahi2015} introduced DeepBind for neural
motif detection \cite{alipanahi2015}.  PhenoSeq uses a statistical approach
combining GC content, 3-mer Shannon entropy, and seven regulatory motif
patterns, with explicit comparison against UCSC phastCons100way scores.

% ═══════════════════════════════════════════════════════════════════════════
\section{Model Description}
% ═══════════════════════════════════════════════════════════════════════════

PhenoSeq consists of two integrated modules: the \emph{HPO Phenotype Extractor}
and the \emph{DNA Conservation Analyser}.

\subsection{Language Detection and Routing}

Input text is analysed by a character-level language detector.  Let
$f_\mathrm{Cyr}$ and $f_\mathrm{Lat}$ denote the fractions of Cyrillic and
Latin characters, respectively.  The routing rule is:
\[
  \text{pipeline} =
  \begin{cases}
    \text{Russian} & \text{if } f_\mathrm{Cyr} > 0.7, \\
    \text{English} & \text{if } f_\mathrm{Lat} > 0.8, \\
    \text{Russian} & \text{otherwise (mixed)}.
  \end{cases}
\]

\subsection{English NER Pipeline}

The English pipeline combines two complementary layers.

\textbf{NER layer.}  \texttt{d4data/biomedical-ner-all} is run with
\texttt{aggregation\_strategy="max"}.  Accepted entity types:
\texttt{Disease\_disorder}, \texttt{Sign\_symptom},
\texttt{Biological\_structure}, \texttt{Diagnostic\_procedure},
\texttt{Lab\_value}.

\textbf{Dictionary layer.}  Regex longest-match over 42{,}355 HPO term names
and synonyms, with span-overlap resolution favouring longer matches.  A
blacklist of 18 clinical modifiers (e.g., \textit{mri, ct, eeg, intermittent,
axial, cerebellar, bilateral, mild, moderate, severe}) prevents procedural and
degree terms from being extracted as phenotypes.

\textbf{Evaluation.}  On a 30-sentence annotated corpus with matching criteria
$\mathrm{Jaccard} \geq 0.6$ and length-ratio $\geq 0.5$:
\[
  P = 0.620, \quad R = 0.838, \quad F_1 = 0.713.
\]
Primary failure mode: multi-word span fragmentation (e.g., ``cerebellar
hypoplasia'' split into ``cerebellar'' + ``hypoplasia''), accounting for 4 of 6
false negatives (see Section~\ref{sec:error}).

\subsection{Russian NER Pipeline}

The Russian pipeline uses a four-layer hybrid architecture (Table~\ref{tab:runerlayers}).

\begin{table}[H]
\centering
\caption{Russian NER pipeline layers.}
\label{tab:runerlayers}
\begin{tabular}{clll}
\toprule
\textbf{Layer} & \textbf{Method} & \textbf{Model / Tool} & \textbf{Confidence} \\
\midrule
1 & HPO Dictionary (RU) & 107-term dict w.\ case forms + longest-match & 0.95 \\
2 & NER Model & \texttt{graviada/labse-ner-runne-ru} & model score \\
3 & Cross-lingual Retrieval & \texttt{paraphrase-multilingual-MiniLM-L12-v2} & $\geq 0.55$ \\
  &                          & (384-dim, 2{,}000-term index) & \\
4 & Rule-based & Cyrillic regex + phenotype prefixes & 0.30--0.50 \\
  &            & + 80 stop-words & \\
\bottomrule
\end{tabular}
\end{table}

\textbf{Layer~1} matches against a hand-curated dictionary inflected across
nominative, accusative, genitive, and instrumental case forms, plus 28
informal/colloquial Russian phrases (e.g.,
\textit{плохо держит голову, мышцы вялые}).
\textbf{Layer~2} provides neural NER coverage via a pre-trained Russian
biomedical NER model.
\textbf{Layer~3} uses a 2{,}000-term semantic index (capped at 500 terms under
low-RAM conditions) for cross-lingual retrieval.
\textbf{Layer~4} catches remaining Cyrillic phrases matching clinical prefixes
(\textit{гипо-}, \textit{гипер-}, \textit{дис-}, etc.).

Key features include 80 clinical stop-words filtering noise (e.g.,
\textit{пациент, МРТ, выполнено}), pymorphy2 lemmatisation (gracefully disabled
on Python~3.11+ due to a known incompatibility), and LaBSE queries capped at
10~per document for CPU latency control.  All dictionaries and indices are
cached with OBO fingerprint validation.

\textbf{Evaluation.}  On an 8-sentence Russian evaluation set:
\[
  P = 0.765, \quad R = 0.765, \quad F_1 = 0.765.
\]

\subsection{HPO Linking}

HPO linking improves on vanilla cosine similarity with three components.

\textbf{Synonym-expanded index.}  The embedding index contains 43{,}596 rows
covering all HPO canonical term names \emph{and} synonyms, built with
\texttt{paraphrase-multilingual-MiniLM-L12-v2} (384-dimensional embeddings).

\textbf{Hierarchy-adjusted re-ranking.}  A depth bonus
\begin{equation}
  \delta(t) = \alpha \cdot \frac{\mathrm{depth}(t)}{\mathrm{depth}_{\max}},
  \quad \alpha = 0.05,
  \label{eq:depth}
\end{equation}
is added to the cosine score of each candidate $t$, favouring more specific
(deeper) ontology nodes and breaking ties in favour of informative terms.

\textbf{Cross-encoder reranker (optional).}
\texttt{cross-encoder/ms-marco-MiniLM-L-6-v2} re-ranks the top-5 bi-encoder
candidates.  Dictionary hits from the NER stage receive a confidence score of
0.98 (not 1.0), allowing the embedding-based ranking to override dictionary
matches when semantically warranted.

\subsection{Orphanet Differential Diagnosis}

Given extracted HPO identifiers, the system queries the Orphanet API for
approximate phenotype matching, returning ranked rare disease candidates.  When
the API is unavailable, a mock dataset covering pontocerebellar hypoplasia
subtypes is provided for demonstration.

\subsection{DNA Conservation Analysis}

The statistical conservation module scores each position using three signals:

\begin{enumerate}[nosep]
  \item \textbf{GC content} $\geq 50\%$: thermodynamic stability proxy
        (weight~1.0).
  \item \textbf{3-mer Shannon entropy}
        $H = -\sum_{k} p_k \log_2 p_k \leq 1.9$:
        low complexity indicates conservation for short motifs (weight~1.0).
  \item \textbf{Regulatory motif hits}: TATA box, GC box (Sp1), CAAT box,
        E-box, Kozak sequence, CpG islands, polyadenylation signal
        (weight~$2.5\times$).
\end{enumerate}

A sliding window of 20~bp (step~1~bp) computes composite scores.  Regions with
normalised score $\geq 0.4$ are labelled \emph{conservative}; gaps $\leq 5$~bp
between conservative regions are merged.

\textbf{Validation against phastCons.}  We compared our statistical scores
against published UCSC phastCons100way scores for the PTEN promoter region
(chr10:89{,}623{,}195--89{,}623{,}395, 200~bp).  Spearman
$\rho = -0.231$ ($p < 0.001$), indicating \emph{weak negative} correlation.
This confirms that sequence-intrinsic signals cannot substitute for
multi-species phylogenetic conservation scores (see Section~\ref{sec:rq3}).

\textbf{Clinical utility.}
PhenoSeq links a patient's symptoms directly to DNA-level analysis through a
three-step workflow.  First, the NER pipeline extracts phenotype terms from the
clinical narrative and maps them to HPO identifiers.  Second, the Orphanet
module identifies rare disease candidates together with their associated genes.
Third, ConservSeq scores the DNA regions within those genes and labels the most
conserved positions.  When a clinician subsequently receives a sequencing report
listing hundreds of variants, they can prioritise variants that fall within
these conserved regions, since a mutation in a functionally important region is
substantially more likely to be causative than one in an arbitrary location.
In this way, PhenoSeq bridges clinical phenotyping and genomic variant
interpretation within a single, end-to-end pipeline.

\subsection{Conservative Sequence Generation}

A GC-optimised greedy generator produces high-GC sequences by selecting at each
position the nucleotide that maximises a weighted score (60\% GC content,
40\% low entropy).  This replaces any ESM2-based approach.  On the demonstration
sequence the generator produces 6~conserved regions (35\% conservative fraction)
and a mean conservation score of 0.771.

% ═══════════════════════════════════════════════════════════════════════════
\section{Evaluation}
% ═══════════════════════════════════════════════════════════════════════════

\subsection{Test Set}

We assembled a 24-case bilingual test set with 47~HPO annotations
(Table~\ref{tab:testset}).

\begin{table}[H]
\centering
\caption{Evaluation test set composition.}
\label{tab:testset}
\begin{tabular}{llrl}
\toprule
\textbf{Group} & \textbf{Cases} & \textbf{Annotations} & \textbf{Description} \\
\midrule
\texttt{en\_formal}   & 8 & 19 & English formal clinical genetics text \\
\texttt{en\_template} & 4 &  5 & English template sentences \\
\texttt{ru\_formal}   & 6 & 17 & Russian formal clinical text \\
\texttt{ru\_informal} & 4 &  6 & Russian informal/colloquial language \\
\texttt{en\_negative} & 1 &  0 & English text with no phenotypes \\
\texttt{ru\_negative} & 1 &  0 & Russian text with no phenotypes \\
\midrule
\textbf{Total}        & \textbf{24} & \textbf{47} & \\
\bottomrule
\end{tabular}
\end{table}

\subsection{Metrics}

\textbf{NER.}  Entity-level Precision ($P$), Recall ($R$), and $F_1$ with
matching criteria Jaccard~$\geq 0.6$ and length-ratio~$\geq 0.5$.

\textbf{HPO Linking.}
\begin{itemize}[nosep]
  \item $\text{Hit@}k$, $k \in \{1,3,5\}$: fraction of queries where the gold
        HPO term appears in the top-$k$ predictions.
  \item MRR: mean reciprocal rank of the gold term,
        $\mathrm{MRR} = \frac{1}{|Q|}\sum_{q} \frac{1}{\mathrm{rank}_q}$.
  \item Wu--Palmer similarity \cite{wu1994}:
        \begin{equation}
          \mathrm{WP}(a,b) =
          \frac{2\cdot\mathrm{depth}\!\left(\mathrm{LCA}(a,b)\right)}
               {\mathrm{depth}(a)+\mathrm{depth}(b)} \;\in [0,1].
          \label{eq:wp}
        \end{equation}
\end{itemize}
All metrics are reported with bootstrap 95\% confidence intervals ($n{=}1{,}000$
resamples).

\subsection{Results}

\subsubsection*{Overall HPO Linking Performance}

Table~\ref{tab:overall} presents per-group and overall results.

\begin{table}[H]
\centering
\caption{Overall HPO linking performance by language/formality group.
95\% bootstrap CIs shown for Hit@1 and MRR where available.}
\label{tab:overall}
\setlength{\tabcolsep}{5pt}
\begin{tabular}{lrcccccc}
\toprule
\textbf{Group} & $n$ & \textbf{Hit@1} & \textbf{Hit@3} & \textbf{Hit@5}
               & \textbf{MRR} & \textbf{WP@1} & \textbf{WP@3} \\
\midrule
\texttt{en\_formal}
  & 19 & 0.368\,{\scriptsize[0.158,\,0.579]}
       & 0.789 & 0.842
       & 0.575\,{\scriptsize[0.404,\,0.733]}
       & 0.368 & 0.842 \\
\texttt{en\_template}
  &  5 & 0.600\,{\scriptsize[0.200,\,1.000]}
       & 1.000 & 1.000
       & 0.800\,{\scriptsize[0.600,\,1.000]}
       & 0.800 & 1.000 \\
\texttt{ru\_formal}
  & 17 & 0.353\,{\scriptsize[0.118,\,0.588]}
       & 1.000 & 1.000
       & 0.627\,{\scriptsize[0.490,\,0.780]}
       & 0.353 & 1.000 \\
\texttt{ru\_informal}
  &  6 & 0.167\,{\scriptsize[0.000,\,0.500]}
       & 0.333 & 0.333
       & 0.250\,{\scriptsize[0.000,\,0.583]}
       & 0.167 & 0.500 \\
\midrule
\textbf{Overall}
  & \textbf{47}
       & \textbf{0.362}\,{\scriptsize[0.234,\,0.489]}
       & \textbf{0.830} & \textbf{0.851}
       & \textbf{0.576}\,{\scriptsize[0.481,\,0.672]}
       & \textbf{0.383} & \textbf{0.872} \\
\bottomrule
\end{tabular}
\end{table}

\subsubsection*{Performance by Symptom Category}

Table~\ref{tab:symptom} breaks down results by phenotype category.

\begin{table}[H]
\centering
\caption{HPO linking performance by symptom group.}
\label{tab:symptom}
\begin{tabular}{lrccc}
\toprule
\textbf{Category} & $n$ & \textbf{Hit@1} & \textbf{Hit@3} & \textbf{MRR} \\
\midrule
Neurological      & 29 & 0.345 & 0.828 & 0.560 \\
Ophthalmological  & 13 & 0.462 & 0.769 & 0.603 \\
Musculoskeletal   &  3 & 0.333 & 1.000 & 0.667 \\
Dysmorphic        &  2 & 0.000 & 1.000 & 0.500 \\
\bottomrule
\end{tabular}
\end{table}

\subsubsection*{Ablation Study (RQ2)}
\label{sec:rq2}

The ablation is conducted on English test cases only (13 cases, 24 annotations),
where the pipeline components can be cleanly toggled (Table~\ref{tab:ablation}).

\begin{table}[H]
\centering
\caption{Ablation study on English test subset (13~cases, 24~annotations).}
\label{tab:ablation}
\begin{tabular}{lccc}
\toprule
\textbf{Configuration} & \textbf{Hit@1} & \textbf{Hit@3}
                       & \textbf{MRR}\,{\scriptsize[95\% CI]} \\
\midrule
Vanilla (cosine only, term-name index)
  & 0.417 & 0.875 & 0.632\,{\scriptsize[0.493,\,0.764]} \\
$+$Synonyms (43{,}596-row index)
  & 0.417 & 0.875 & 0.635\,{\scriptsize[0.494,\,0.778]} \\
$+$Hierarchy (depth bonus $\alpha{=}0.05$)
  & 0.417 & 0.875 & 0.632\,{\scriptsize[0.493,\,0.785]} \\
$+$Cross-encoder (\texttt{ms-marco-MiniLM-L-6-v2})
  & 0.417 & 0.792 & 0.614\,{\scriptsize[0.474,\,0.751]} \\
\bottomrule
\end{tabular}
\end{table}

On this formal English subset, where all test terms possess exact string matches
in the HPO synonym index, the vanilla baseline already performs well
($\mathrm{MRR} = 0.632$).  The synonym-expanded index provides a marginal gain
of $+0.003$ MRR ($+0.5\%$ relative).  Hierarchy bonus yields no additional
improvement.  Cross-encoder reranking \emph{reduces} performance
($\Delta\text{Hit@3} = -0.083$, $\Delta\text{MRR} = -0.018$), attributable to
domain mismatch between the MS~MARCO training corpus and biomedical text.
These additional components therefore serve as insurance for harder cases
(informal language, non-canonical synonyms) rather than improving metrics on
the current formal test set.

\subsubsection*{DNA Conservation (RQ3)}
\label{sec:rq3}

Spearman $\rho = -0.231$ ($p < 0.001$) between PhenoSeq statistical scores
and phastCons100way scores on the PTEN promoter region
(chr10:89{,}623{,}195--89{,}623{,}395).  The weak \emph{negative} correlation
confirms that GC content, Shannon entropy, and regulatory motif signals are
insufficient proxies for multi-species phylogenetic conservation.  We report
this as a transparent negative result: sequence-intrinsic statistics cannot
replace multi-species alignment-based conservation scoring.

% ═══════════════════════════════════════════════════════════════════════════
\section{Error Analysis}
\label{sec:error}
% ═══════════════════════════════════════════════════════════════════════════

\subsection{Primary Failure Mode: Incomplete Mention Spans}

The English NER model performs token-level classification and struggles with
multi-word phenotype boundaries.  Terms such as ``cerebellar hypoplasia'' and
``flexion contracture'' are split into constituent tokens, causing the
downstream linking step to receive incomplete input.  This accounts for 4 of
6~false negatives on the 30-sentence annotated corpus.

\subsection{Cross-encoder Domain Mismatch}

\texttt{cross-encoder/ms-marco-MiniLM-L-6-v2}, trained on general passage
ranking (MS~MARCO), reranked correct biomedical matches below lexically similar
but ontologically incorrect candidates.  This caused $\text{Hit@3}$ to drop
from 0.875 to 0.792 on formal English.  A biomedical cross-encoder fine-tuned
on HPO term pairs would be required for reliable reranking.

\subsection{Russian Informal Language}

Russian informal and colloquial clinical descriptions (\textit{глаз косит},
\textit{плохо держит голову}) achieve $\text{Hit@3} = 0.333$, compared to
1.000 for formal Russian.  The 107-term fallback dictionary augmented with
28~informal phrases provides coverage for the most common lay descriptions but
is insufficient for the full range of colloquial expression.

\subsection{Ranking Gap}

A 47-percentage-point gap between $\text{Hit@1} = 0.362$ and
$\text{Hit@3} = 0.830$ indicates that the correct HPO term is consistently
retrieved within the top-3 candidates but frequently not ranked first.  This
is a \emph{ranking} problem rather than a retrieval problem, and is the primary
direction for future improvement.

\subsection{Systematic Error Patterns}

\begin{itemize}[nosep]
  \item \textbf{Multi-word boundaries}: ``pontine hypoplasia'' split to
        ``pontine'' $+$ ``hypoplasia''.
  \item \textbf{Non-canonical synonyms}: descriptions absent from the HPO
        synonym list require expanded dictionary coverage.
  \item \textbf{Genitive case (Russian)}: morphological variants not fully
        covered by the 107-term dictionary.
  \item \textbf{Negation}: ``No seizures'' would be incorrectly extracted as
        ``seizures''; negation detection is not yet implemented.
\end{itemize}

% ═══════════════════════════════════════════════════════════════════════════
\section{Reproducibility}
% ═══════════════════════════════════════════════════════════════════════════

\begin{itemize}[nosep]
  \item Random seed fixed (\texttt{SEED = 42}).
  \item HPO OBO file cached locally with MD5 fingerprint validation;
        cache invalidates automatically on HPO updates.
  \item All embedding indices cached with OBO fingerprint.
  \item All models run on CPU (no GPU required).
  \item \texttt{requirements.txt} provided with pinned versions.
  \item Full test suite: \textbf{29/30 tests passed (96\%)}.
\end{itemize}

\textbf{Key dependencies:} \texttt{torch}~2.12.0, \texttt{transformers}~5.8.1,
\texttt{sentence-transformers}~5.5.1, \texttt{numpy}~1.26.4,
\texttt{pandas}~2.2.2, \texttt{scikit-learn}~1.8.0, \texttt{networkx}~3.3.

% ═══════════════════════════════════════════════════════════════════════════
\section{Limitations}
% ═══════════════════════════════════════════════════════════════════════════

\begin{description}[nosep]
  \item[NER training.] English NER uses a general biomedical model, not
        HPO-specific.  Fine-tuning on PhenoTagger annotations would improve
        $F_1$.
  \item[Russian NER evaluation.] Only 8~annotated sentences; a larger test set
        is needed for statistical significance.
  \item[DNA conservation.] Statistical signals show weak negative correlation
        with phastCons ($\rho = -0.231$).  Multi-species alignment is required
        for publishable conservation analysis.
  \item[Test set size.] 24~cases / 47~annotations is sufficient for a student
        project but not for clinical validation ($\geq 200$ annotated records
        required).
  \item[Negation handling.] Not implemented; ``no seizures'' would be
        incorrectly extracted as a positive phenotype.
  \item[pymorphy2.] Non-functional on Python~3.11+; Russian lemmatisation is
        gracefully disabled.
  \item[Cross-encoder.] The current cross-encoder degrades performance on
        biomedical text due to domain mismatch; a biomedical-specific variant
        is required.
\end{description}

% ═══════════════════════════════════════════════════════════════════════════
\section{Conclusion}
% ═══════════════════════════════════════════════════════════════════════════

PhenoSeq demonstrates that a combined NER + semantic linking pipeline can
achieve clinically useful phenotype extraction in both English and Russian using
only CPU-compatible models.  On formal clinical genetics text, $\text{Hit@3}$
reaches 0.789--1.000 across languages.  Russian informal language remains the
principal challenge ($\text{Hit@3} = 0.333$).  The DNA conservation module
provides interpretable statistical signals but must be considered a
demonstration, not a replacement for phylogenetic conservation scores
($\rho = -0.231$ vs.\ phastCons).

\textbf{Key findings:}
\begin{itemize}[nosep]
  \item English NER: $F_1 = 0.713$ (30-sentence annotated corpus,
        $P = 0.620$, $R = 0.838$).
  \item Russian NER: $F_1 = 0.765$ (8-sentence evaluation set).
  \item HPO linking: $\text{Hit@3} = 0.830$ overall; Russian formal
        $= 1.000$; Russian informal $= 0.333$.
  \item A 47-point Hit@1--Hit@3 gap ($0.362 \to 0.830$) identifies ranking
        as the primary bottleneck.
  \item Synonym expansion provides marginal MRR gain ($+0.003$);
        cross-encoder reranking decreases performance on the current test set
        ($\Delta\text{Hit@3} = -0.083$) due to domain mismatch.
  \item Statistical DNA conservation (ConservSeq): a lightweight, interpretable
        method for identifying functionally important genomic regions.  The
        module bridges clinical phenotyping and genomic variant interpretation
        --- given a set of symptoms, it identifies candidate genes and their
        most conserved regions, helping clinicians prioritise which mutations
        to investigate first in sequencing data.  Validated against phastCons
        (Spearman $\rho = -0.231$).
\end{itemize}

\textbf{Future work:}
\begin{itemize}[nosep]
  \item Collecting annotated Russian clinical referral letters to address the
        informal language gap.
  \item Fine-tuning RuBERT on NEREL-BIO for Russian NER.
  \item Integrating phyloP/GERP scores via Ensembl REST API for publishable
        conservation analysis.
  \item Clinical validation on MIMIC-III discharge summaries.
  \item Negation detection using NegEx or a spaCy dependency parser.
  \item Replacing the cross-encoder with a biomedical variant fine-tuned on
        HPO term pairs.
\end{itemize}

% ═══════════════════════════════════════════════════════════════════════════
\bibliographystyle{plain}
\begin{thebibliography}{99}

\bibitem{kohler2021}
K\"ohler S et al.\ (2021).
The Human Phenotype Ontology in 2021.
\textit{Nucleic Acids Research}, 49(D1):D1207--D1217.

\bibitem{robinson2010}
Robinson PN, Mundlos S (2010).
The Human Phenotype Ontology.
\textit{Clinical Genetics}, 77(6):525--534.

\bibitem{jacobsen2022}
Jacobsen JOB et al.\ (2022).
The GA4GH Phenopackets standard.
\textit{Nature Genomics}, 4(4):44.

\bibitem{groza2015}
Groza T et al.\ (2015).
Automatic concept recognition using the HPO reference and test suite corpora.
\textit{Database}, 2015:bav022.

\bibitem{nevéol2018}
N\'ev\'eol A et al.\ (2018).
Clinical NLP for languages other than English.
\textit{JAMIA}, 25(12):1449--1467.

\bibitem{lee2020}
Lee J et al.\ (2020).
BioBERT: a pre-trained biomedical language representation model.
\textit{Bioinformatics}, 36(4):1234--1240.

\bibitem{huang2022}
Huang L et al.\ (2022).
PhenoBERT: Combining individual models for HPO term identification.
\textit{IEEE/ACM TCBB}, 20(1):943--952.

\bibitem{neumann2019}
Neumann M et al.\ (2019).
ScispaCy: Fast and robust models for biomedical NLP.
\textit{18th BioNLP Workshop}, 319--327.

\bibitem{miftah2017}
Miftahutdinov Z et al.\ (2017).
Identifying disease-related expressions in news articles.
\textit{Dialog conference}, 1:284--293.

\bibitem{louka2021}
Loukachevitch N et al.\ (2021).
NEREL-BIO: a dataset of biomedical abstracts.
\textit{RANLP 2021}, 413--422.

\bibitem{xiong2021}
Xiong L et al.\ (2021).
ANCE: Approximate nearest neighbour negative contrastive estimation.
\textit{ICLR 2021}.

\bibitem{nogueira2020}
Nogueira R, Cho K (2020).
Passage re-ranking with BERT.
\textit{arXiv}:1901.04085.

\bibitem{kohler2009}
K\"ohler S et al.\ (2009).
Clinical diagnostics in human genetics with semantic similarity searches.
\textit{AJHG}, 85(4):457--464.

\bibitem{pollard2010}
Pollard KS et al.\ (2010).
Detection of nonneutral substitution rates on mammalian phylogenies.
\textit{Genome Research}, 20(1):110--121.

\bibitem{alipanahi2015}
Alipanahi B et al.\ (2015).
Predicting sequence specificities of DNA- and RNA-binding proteins.
\textit{Nature Biotechnology}, 33(8):831--838.

\bibitem{robinson2020}
Robinson PN et al.\ (2020).
LIRICAL: likelihood ratio interpretation of clinical abnormalities.
\textit{AJHG}, 107(3):403--417.

\bibitem{nguengang2020}
Nguengang Wakap S et al.\ (2020).
Estimating cumulative point prevalence of rare diseases.
\textit{EJHG}, 28(2):165--173.

\bibitem{wu1994}
Wu Z, Palmer M (1994).
Verbs semantics and lexical selection.
\textit{ACL 1994}, 133--138.

\end{thebibliography}

% ═══════════════════════════════════════════════════════════════════════════
\appendix
% ═══════════════════════════════════════════════════════════════════════════

\section{Module Test Suite Results}

29/30 tests passed (96\%).  One reproducibility test failed (see below).

\begin{table}[H]
\centering
\caption{Full test suite results.}
\label{tab:tests}
\begin{tabular}{lrl}
\toprule
\textbf{Category} & \textbf{Result} & \textbf{Tests included} \\
\midrule
NER           & 9/9  & No \texttt{\#\#} tokens, dict layer, RU stop-word \\
              &      & filtering, multi-word extraction, score bounds \\
Wu--Palmer    & 3/3  & Range validation, identity check \\
HPO Linking   & 4/4  & Results returned, correct HPO found, \\
              &      & hierarchy bonus applied, synonym matching \\
Orphanet      & 2/2  & API call, mock fallback \\
DNA Conservation & 4/4 & GC signal, entropy signal, motif detection, \\
              &      & region merging \\
Evaluation    & 5/5  & Test set size, annotation count, negative \\
              &      & controls, bootstrap CI, results populated \\
Reproducibility & 2/3 & \texttt{requirements.txt}, HPO fingerprint \\
              &      & (\textit{1 failure}: embedding cache fingerprint) \\
\midrule
\textbf{Total} & \textbf{29/30 (96\%)} & \\
\bottomrule
\end{tabular}
\end{table}

\section{Sample Outputs}

\subsubsection*{English Formal Input}

\textit{``The patient presents with cerebellar hypoplasia, axial hypotonia,
ptosis, and intermittent seizures.  MRI shows pontine hypoplasia.''}

\medskip
\noindent Extracted phenotypes:
\begin{itemize}[nosep]
  \item \textit{pontine hypoplasia} $\to$ \texttt{HP:0012110} (score~0.98)
  \item \textit{seizures} $\to$ \texttt{HP:0001250} (score~0.98)
  \item \textit{ptosis} $\to$ \texttt{HP:0000508} (score~0.98)
\end{itemize}

\noindent Orphanet differential:
ORPHA:2524 Pontocerebellar hypoplasia type~1A (71\%),
ORPHA:99803 Pontocerebellar hypoplasia type~2A (65\%).

\subsubsection*{Russian Formal Input}

\textit{``Пациент 4~лет.  Аксиальная гипотония, птоз левого века, косоглазие.
МРТ: гипоплазия мозжечка.''}

\medskip
\noindent Extracted phenotypes:
\begin{itemize}[nosep]
  \item \textit{гипоплазия мозжечка} $\to$ \texttt{HP:0001321} (score~0.95)
  \item \textit{аксиальная гипотония} $\to$ \texttt{HP:0008936} (score~0.95)
  \item \textit{косоглазие} $\to$ \texttt{HP:0000486} (score~0.95)
  \item \textit{птоз} $\to$ \texttt{HP:0000508} (score~0.95)
\end{itemize}

\noindent Orphanet differential:
ORPHA:2524 Pontocerebellar hypoplasia type~1A (71\%).

\subsubsection*{Russian Informal Input}

\textit{``Ребёнок плохо держит голову, мышцы очень вялые, судороги.''}

\medskip
\noindent Extracted phenotypes:
\begin{itemize}[nosep]
  \item \textit{плохо держит голову} $\to$ \texttt{HP:0008936} (score~0.95)
  \item \textit{судороги} $\to$ \texttt{HP:0001250} (score~0.95)
  \item \textit{мышцы очень вялые} $\to$ \texttt{HP:0001324} (score~0.81)
\end{itemize}

\end{document}
